
\section{模型改进与推广}

尽管本文所建立的模型在【研究问题】中均取得了合理的结果,但仍存在一些可改进之处,值得在后续研究中进一步探索。

在数据层面,本文所使用的数据集【数据局限说明】。后续研究应考虑纳入更广泛的数据,以提高模型的普适性。

在方法层面,【说明模型局限,提出改进方向】。后续研究可以引入【新方法】以处理【现有方法不足】。

在临床应用层面,本文所给出的【推荐方案】是基于【特定假设】得到的,【参数/权重】的取值具有【主观性】。后续研究可以通过【临床调查】获取专家共识,确定更符合实际应用的设置。

在工程实现层面,本文的模型可通过部署为 Web 服务或集成至【目标系统】,为【用户角色】提供实时的【功能】。【主方法】由于其【优势】,尤其适合在【应用场景】中应用。

最后,本文的研究框架可推广至【其他应用场景】,通过扩展特征集合和重新训练模型,实现方法在不同【任务】中的迁移。
